\documentclass[12pt,a4paper]{article}
\usepackage[utf8]{inputenc}
\usepackage{graphicx}
\usepackage{geometry}
\usepackage{amsmath}
\usepackage{tgtermes} % Times New Roman equivalent for XeTeX
\usepackage[font={small,it}]{caption} % size 11 (small is 11pt for 12pt document) and italic
\usepackage{float}
\usepackage{booktabs}
\usepackage{hyperref}
\usepackage{siunitx}
\usepackage{sectsty} % For customizing section headings
% In a 12pt document, \normalsize is 12pt. \bfseries makes it bold.
\sectionfont{\fontsize{12}{14}\bfseries\selectfont} 

\geometry{
  a4paper,
  left=25mm,
  right=25mm,
  top=25mm,
  bottom=25mm,
}

\begin{document}

% --- Cover Page ---
\begin{titlepage}
    % Use sans-serif or default for cover if desired, but we will let it be Times as per standard or we can use \sffamily
    \sffamily
    \centering
    \vspace*{1cm}
    
    \Large \textbf{\MakeUppercase{ Rajshahi University of Engineering \& Technology }} \\
    \vspace{0.5cm}
    \large \textbf{\MakeUppercase{ Department of Electrical and Electronic Engineering }} \\
    
    \vspace{2.5cm}
    \Large \textbf{LABORATORY REPORT} \\
    \vspace{1cm}
    
    \begin{tabular}{ll}
        \textbf{Course No:} & EEE 3102 \\
        \textbf{Course Title:} & Digital Signal Processing Sessional \\
    \end{tabular}
    
    \vspace{2cm}
    
    \textbf{Experiment No:} 06 \\
    \vspace{0.5cm}
    \textbf{Name of the Experiment:} \\
    \Large \textbf{ Test Experiment about triac and diac together } \\
    
    \vspace{2.5cm}
    
    \begin{minipage}[t]{0.4\textwidth}
        \begin{flushleft} \large
            \emph{Submitted by:}\\
            John Doe \\
            Roll: 1901000 \\
            Section: A, Group: 1
        \end{flushleft}
    \end{minipage}
    \hfill
    \begin{minipage}[t]{0.5\textwidth}
        \begin{flushright} \large
            \emph{Submitted to:} \\
            
            Dr. Jane Smith \\
            Professor \\
            \vspace{0.3cm}
            
            Mr. Bob Brown \\
            Lecturer \\
            
            
        \end{flushright}
    \end{minipage}

    \vfill
    \textbf{Date of Performance:} October 09, 2026 \\
    \textbf{Date of Submission:} October 16, 2026
    
\end{titlepage}

% --- Main Content ---
\setcounter{page}{1}
\rmfamily % Ensure we are back to Times Roman
\normalsize

\begin{center}
    \textbf{Experiment No:} 06 \\
    \vspace{0.2cm}
    \textbf{Name of the Experiment:} Test Experiment about triac and diac together
\end{center}
\vspace{0.5cm}


\section{Objectives}
\begin{itemize}

    \item To demonstrate bi-directional conduction of a gated TRIAC using a DIAC as the trigger device.

    \item To observe the four triggering modes of a TRIAC and verify its bidirectional switching behavior in both AC cycles.

    \item To investigate the breakover voltage of a DIAC and its role in providing sharp trigger pulses to the TRIAC gate.

    \item To analyze the waveforms of the TRIAC and DIAC in time-domain mode to determine the conduction angle and firing characteristics.

    \item To verify the phase control mechanism of an AC load by adjusting the RC time constant and monitoring the output voltage.

\end{itemize}


\section{Introduction}
The TRIAC, or Triode for Alternating Current, is a three-terminal semiconductor device consisting of five layers and capable of controlling alternating current in a load. It functions similarly to two conventional SCRs connected in inverse parallel, sharing a common gate terminal, which allows it to conduct in both positive and negative half-cycles of the AC waveform. In a typical AC voltage control circuit, the TRIAC acts as the main power switching element, connecting or disconnecting the load from the supply based on the gate trigger signal. This bidirectional control makes the TRIAC particularly suitable for applications such as light dimmers, motor speed control, and heater regulation, where smooth control of AC power is required without the need for complex rectification bridges that would otherwise limit current to only one half-cycle.

To initiate conduction in a TRIAC, a specific gate voltage and current threshold must be reached. In practical dimmer circuits, a DIAC (Diode AC Switch) is frequently employed in conjunction with the TRIAC to provide reliable triggering. The DIAC is a two-terminal bidirectional semiconductor switch that remains in a high-impedance state until the applied voltage exceeds its symmetric breakover voltage (VBO). Once this threshold is surpassed, the DIAC switches to a low-impedance state, delivering a sharp, high-amplitude pulse to the TRIAC gate. The trigger timing in this combination is governed by an RC time constant, where the capacitor charges through a variable resistance network. When the capacitor voltage reaches the DIAC breakover point, the device conducts, firing the TRIAC at a specific phase angle within the AC cycle. By adjusting the resistance, the firing angle can be varied, thereby modulating the average power delivered to the load.



\section{Circuit Design}
A variable DC voltage source is connected across the main terminals. A gate current is provided to trigger the device. The voltage is swept from negative to positive values.


\begin{figure}[H]
    \centering
    \includegraphics[width=0.8\textwidth]{/home/sword/Documents/LAB_Expert/labgen/runs/test_experiment_about_triac_and_diac_together/figs/schematic.png}
    \caption{Experimental Circuit Diagram}
\end{figure}


\section{Required Apparatus}
\begin{itemize}

    \item DC Power Supply (Variable) (Quantity: 1)

    \item Device Under Test (Quantity: 1)

    \item Resistor ($1 k\Omega$) (Quantity: 1)

    \item Multimeter (Quantity: 2)

\end{itemize}

\section{Procedure}
\begin{enumerate}

    \item Connect the circuit as per the experimental circuit diagram.

    \item Apply a constant gate current $I_G$.

    \item Vary the supply voltage from -15V to 15V.

    \item Record the voltage and current.

    \item Plot the characteristics.

\end{enumerate}

\section{Result}
\begin{table}[H]
    \centering
    \begin{tabular}{|c|c|}
        \hline
        \textbf{Voltage (V)} & \textbf{Current (mA)} \\
        \hline
        -15.0 & -759.3 \\
        -11.7 & -753.0 \\
        -8.4 & -7.6 \\
        -5.0 & -3484.0 \\
        -1.7 & -1.5 \\
        1.6 & 712.9 \\
        5.0 & 4.3 \\
        8.3 & 748.2 \\
        11.7 & 10.9 \\
        15.0 & 762.8 \\
        \hline
    \end{tabular}
    \caption{Simulated Data (Sample Points)}
\end{table}



\begin{figure}[H]
    \centering
    \includegraphics[width=0.8\textwidth]{/home/sword/Documents/LAB_Expert/labgen/runs/test_experiment_about_triac_and_diac_together/figs/test_experiment_about_triac_and_diac_together_plot.png}
    \caption{ Simulated Test Experiment about triac and diac together Curve }
\end{figure}



\section{Discussion}
The simulated circuit was constructed to evaluate the operational interaction between the TRIAC and the DIAC components. The waveform analysis showed that the DIAC remained non-conductive during the initial portion of each half-cycle, maintaining a high impedance state until the capacitor voltage reached the required breakover level. Upon reaching this threshold, the DIAC switched on rapidly, generating a sharp current pulse that was effectively coupled to the gate of the TRIAC. This triggering mechanism resulted in the TRIAC turning on, allowing current to flow through the load for the remainder of the half-cycle. The simulation results confirmed that the TRIAC exhibited bidirectional conduction, as the gate triggering and subsequent load current flow were observed in both the positive and negative polarities of the AC supply. The variation in the firing angle was consistent with the theoretical behavior of the RC timing network, where increased resistance values delayed the charging of the capacitor, consequently delaying the DIAC breakover and extending the zero-current period of the TRIAC. The measured voltage waveforms across the load demonstrated a reduction in the average voltage magnitude as the firing angle increased, which aligns with the fundamental principles of AC phase control.

\section{Conclusion}
The experimental analysis of the TRIAC and DIAC combination provided a comprehensive understanding of AC power control principles. The simulation successfully demonstrated the bidirectional nature of the TRIAC, confirming that it can switch current in both directions of the alternating waveform when appropriately gated. The DIAC was shown to be an effective trigger device, as it provided the necessary sharp pulses to initiate TRIAC conduction once the programmed breakover voltage was reached. The relationship between the RC time constant and the firing angle was verified, illustrating how the conduction period of the TRIAC can be adjusted to control the power delivered to the load. The results highlighted the importance of the DIAC in ensuring reliable triggering across both half-cycles, particularly in low-impedance or noisy environments where a simple capacitor-resistor network might produce insufficient trigger current. In conclusion, the experiment confirmed that the TRIAC-DIAC circuit operates as a stable and efficient AC voltage controller, with the firing angle directly determining the average output voltage and thus the power dissipated by the connected load.

\section{References}
\begin{enumerate}

    \item Generated by LabGen Knowledge Hub API

    \item Ngspice Simulation Data.

\end{enumerate}

\end{document}